\section{Motivation for the Research} \label{section:motivation}

\label{para:nonideal_interface_formation} % ---------- %

Interfaces rarely form under ideal conditions.

Instead, they emerge through complex, often non-equilibrium processes such as high-temperature annealing, epitaxial
growth, and chemically driven surface reconstructions.

These dynamics frequently trap systems in metastable configurations or generate local disorder, resulting in atomic
arrangements that deviate significantly from those predicted by ideal lattice-matching rules.

Such complexity challenges efforts to model interface behaviour using conventional approaches alone.

\label{para:dft_scaling_limitations} % ---------- %

While DFT provides a reliable route to total energy prediction and structural relaxation,
its unfavourable scaling with system size ($\mathcal{O}(N^3)$ in electron count) severely limits its use in
high-throughput workflows.

In particular, DFT becomes prohibitive when applied to interface structures exceeding a few hundred atoms, as required
to capture long‐range strain, multi‐layer reconstruction, or compositional variation.

Moreover, standard DFT approximations, such as semi‐local exchange–correlation functionals, often struggle to predict key
interfacial phenomena in layered systems.

For example, accurate band‐alignment at a heterojunction determines charge‐injection barriers and device turn‐on
voltages; polarization effects at oxide or ferroelectric interfaces give rise to internal electric fields and
screening that strongly modulate carrier densities; and weak van der Waals interactions between two‐dimensional
layers govern adhesion, interlayer spacing, and moiré‐driven electronic reconstructions.

Failure to capture these effects can lead to qualitatively incorrect predictions of interface stability and function.

MLP-based approaches can in principle overcome many of these challenges by incorporating
long‐range interactions and nonlocal correlations learned from data.

However, since MLPs are trained on DFT outputs, their accuracy is fundamentally limited by the fidelity of the
underlying DFT dataset: any systematic errors in band alignments, polarization responses, or dispersion forces will be
inherited by the potential.

Thus, while MLPs enable vastly more efficient sampling of configurational space, their predictive power remains
contingent on high‐quality reference data.

\label{para:mlp_overview} % ---------- %

Machine‐learned interatomic potentials (MLPs) offer a promising route to overcoming these limitations.

By learning to emulate DFT‐level energies and forces, MLPs like MACE provide a means to conduct large‐scale structure
relaxations and energy evaluations at dramatically reduced cost. This enables exploration of broader configurational
spaces, including low‐symmetry systems and extended defects, that would otherwise remain intractable.

However, this advantage is only realised if the models generalise well across diverse interfacial environments.

MLPs offer a promising route to overcoming these limitations.

This capacity opens up the exploration of vastly broader configurational spaces—including low-symmetry interfaces and
extended defects—that would be intractable with DFT alone.

Crucially, however, achieving DFT-comparable accuracy across diverse interfacial environments demands a representative,
high-fidelity training set: any gaps or systematic biases in the reference data will be reflected in the potential’s
predictions.


\label{para:motivation_mlp_dft} % ---------- %

The primary aim of this report is to evaluate whether MLPs can reproduce the DFT‐derived ranking of interface
favourability, here quantified by formation energy, across a set of group 14 heterostructures.

While the overarching PhD objective is to develop a Machine‐Learned Interface (MLI) model that predicts interface
energies directly from the constituent slabs (bypassing explicit interface construction), this study focuses on rapid
MLP‐assisted screening to identify the most promising candidates for eventual full DFT ISIF3 relaxations.

Chapter~\ref{chapter:computational_investigation} shows that MLPs often preserve relative DFT energy rankings,
particularly within chemically similar systems, with a few edge cases like pure‐Si interfaces (grain boundaries).

These discrepancies underscore the need for targeted refinement, motivating the delta‐learning approaches presented in
Chapter~\ref{chapter:next_steps}, which correct systematic MLP–DFT residuals via supervised learning on
interface‐specific error distributions.

\label{para:feature_based_interface_prediction} % ---------- %

In addition to correcting prediction accuracy, this research investigates whether interface favourability can be
inferred directly from structural or surface-level features.

Chapter~\ref{chapter:next_steps} proposes the use of learned models that map surface properties to interfacial energies,
potentially bypassing the need for full supercell construction.

This complements the structure-generation pipelines, based on ARTEMIS, that systematically enumerate viable interface
candidates across various stacking shifts and Miller plane pairings.

\label{para:project_motivation} % ---------- %

Ultimately, this project is motivated by the need for more efficient and predictive tools to accelerate the design of
functional heterostructures.

By embedding MLPs within automated screening workflows and augmenting them with delta-learning corrections and
structure-based heuristics, it becomes possible to bridge the gap between physical realism and computational
tractability.

In doing so, this research contributes to the broader aim of enabling scalable, data-driven exploration of complex
interfacial systems; particularly in regimes where conventional symmetry-based rules or bulk-derived heuristics break
down.